
\chapter{مقدمه}
\label{ch:intro}

گفت‌وگوی گفتاری، طبیعی‌ترین شکل تعامل انسان با ماشین است. با این حال، تا همین
اواخر سامانه‌های گفت‌وگوی گفتاری تقریباً همیشه از زنجیره‌ای از مؤلفه‌های مستقل
ساخته می‌شدند: تشخیص فعالیت گفتاری، بازشناسی گفتار، یک موتور گفت‌وگوی متنی و
سپس تبدیل متن به گفتار. معرفی مدل‌هایی مانند \en{Moshi}~\cite{defossez2024moshi}
و \en{Freeze-Omni}~\cite{wang2025freezeomni} نشان داد که می‌توان گفت‌وگوی گفتاری
را به‌صورت یک مسئله «گفتار‌به‌گفتار» و در قالب یک مدل یکپارچه مدل کرد؛ مدلی که
هم‌زمان می‌شنود و سخن می‌گوید و برای نوبت‌گیری به بخش‌بندی صریح گفتار نیاز
ندارد.

این تحول برای زبان‌های پرمنبع مانند انگلیسی و چینی رخ داده است. برای زبان فارسی
هیچ‌یک از مدل‌های گفتار‌به‌گفتار متن‌باز موجود پشتیبانی رسمی ارائه نمی‌کنند، و
مجموعه داده‌های مکالمه‌ای فارسی که هم‌پوشانی گفتار، وقفه و بازخورد کلامی را
پوشش دهند نیز به‌صورت عمومی در دسترس نیستند. پرسش عملی این پروژه از همین شکاف
برمی‌آید: با بودجه محاسباتی و زمانی یک پروژه کارشناسی، تا چه اندازه می‌توان یک
سامانه گفتار‌به‌گفتار فارسی تمام‌دوطرفه ساخت، و مهم‌تر از آن، چه چیزی را
می‌توان درباره آن با شواهد قابل ممیزی «اثبات» کرد؟

\section{تعریف مسئله}
\label{sec:problem}

هدف تعریف‌شده این پروژه، طراحی، پیاده‌سازی و ارزیابی یک نمونه اولیه از یک مدل
زبانی بزرگ گفتار‌به‌گفتار برای زبان فارسی است. سامانه باید به‌صورت
«تمام‌دوطرفه»\پاورقی{full-duplex} کار کند؛ یعنی برخلاف سامانه‌های
نیم‌دوطرفه که تنها پس از پایان‌یافتن پاسخ به کاربر گوش می‌دهند، میکروفون کاربر
در طول پخش پاسخ نیز فعال بماند و سامانه بتواند در صورت
«وقفه»\پاورقی{interruption} کاربر واکنش نشان دهد.

سه ویژگی این مسئله را از یک سامانه پرسش‌و‌پاسخ گفتاری متعارف جدا می‌کند:

\begin{itemize}
\item \مهم{هم‌زمانی شنیدن و گفتن.} سامانه باید در حالی که خروجی صوتی خود را
      پخش می‌کند، جریان ورودی میکروفون را نیز پردازش کند و درباره ادامه یا
      توقف پخش تصمیم بگیرد.
\item \مهم{تصمیم‌گیری زمان‌حساس.} تصمیم توقف پخش باید در مقیاس چند‌ده
      میلی‌ثانیه گرفته شود، وگرنه کاربر احساس می‌کند سامانه حرف او را
      نمی‌شنود.
\item \مهم{تداوم گفتار قطع‌شده.} گفتاری که کاربر در لحظه وقفه ادا می‌کند
      نباید از دست برود؛ این گفتار باید به نوبت بعدی منتقل شود، در غیر این صورت
      کاربر ناچار می‌شود جمله خود را تکرار کند.
\end{itemize}

\section{اهمیت موضوع}
\label{sec:motivation}

پرداختن به این مسئله برای زبان فارسی از چند جهت اهمیت دارد.

نخست، \مهم{دسترس‌پذیری}. رابط گفتاری برای کاربرانی که با صفحه‌کلید فارسی راحت
نیستند، برای کاربران کم‌بینا و به‌طور خاص برای سالمندان، می‌تواند تنها راه
مؤثر استفاده از خدمات دیجیتال باشد. در گفت‌وگوی طبیعی سالمندان، مکث‌های بلندتر
و اصلاح‌های میان‌جمله‌ای بیشتر رخ می‌دهد؛ یک سامانه نیم‌دوطرفه که تا پایان پاسخ
خود گوش نمی‌دهد، در این شرایط تجربه‌ای آزارنده می‌سازد.

دوم، \مهم{کم‌منبع‌بودن فارسی در این حوزه خاص}. برای فارسی مجموعه داده‌های
بازشناسی گفتار و تبدیل متن به گفتار در دسترس است؛ نمونه شاخص آن
\en{ManaTTS}~\cite{qharabagh2025manatts} با بیش از \N{86} ساعت گفتار تک‌گوینده
است. اما این منابع، تک‌گوینده و خوانشی‌اند و پدیده‌های تعاملی‌ای که یک مدل
تمام‌دوطرفه باید یاد بگیرد، یعنی هم‌پوشانی و وقفه و بازخورد کلامی، در آن‌ها وجود
ندارد.

سوم، \مهم{هزینه محاسباتی}. مدل‌های گفتار‌به‌گفتار روز، مانند
\en{Qwen2.5-Omni}~\cite{xu2025qwen25omni}، برای اجرا به حافظه‌ای بسیار بیش از
یک کارت گرافیک آزمایشگاهی نیاز دارند. نشان‌دادن این‌که چه چیزی در بازه حافظه
\N{12} تا \N{24} گیگابایت «قابل دستیابی» و چه چیزی «قابل دستیابی نیست»،
خودش یک نتیجه مهندسی سودمند است.

\section{نیازمندی‌های تعریف‌شده}
\label{sec:requirements}

تعریف پروژه هفت نیازمندی مشخص را تعیین کرده است. این نیازمندی‌ها در سراسر این
نوشتار به‌عنوان معیار سنجش دستاورد به کار می‌روند و در
جدول~\ref{tab:requirements} فهرست شده‌اند. برای هر نیازمندی، فصل مربوط به
نتایج (فصل~\ref{ch:results}) وضعیت نهایی و کلاس شواهد پشتیبان آن را گزارش
می‌کند.

\begin{table}[ht]
\centering
\caption{نیازمندی‌های هفت‌گانه تعریف پروژه و محل بررسی آن‌ها در این نوشتار}
\label{tab:requirements}
\small
\begin{tabular}{clL{6.2cm}}
\toprule
\textbf{شماره} & \textbf{نیازمندی} & \textbf{سنجه تعیین‌شده} \\
\midrule
ن۱ & نمونه اولیه گفتار‌به‌گفتار فارسی & تولید پاسخ گفتاری فارسی از ورودی
      گفتاری فارسی \\
ن۲ & کارکرد تمام‌دوطرفه & شنیدن در هنگام پخش و واکنش به وقفه \\
ن۳ & تأخیر سرتاسری & حداکثر \N{500} میلی‌ثانیه از دریافت گفتار تا تولید پاسخ
      گفتاری \\
ن۴ & مجموعه داده مکالمه فارسی & \N{100} تا \N{200} ساعت گفتار همراه با نویز،
      هم‌پوشانی و قطع‌و‌وصل کلام \\
ن۵ & ماژول تشخیص وقفه & دقت بالای \N{80} درصد با ویژگی‌های انرژی، زیر‌و‌بمی و
      \en{MFCC} \\
ن۶ & ارزیابی روی سخت‌افزار هدف & پردازنده گرافیکی با \N{12} تا \N{24} گیگابایت
      حافظه \\
ن۷ & ارزیابی انسانی & \N{5} تا \N{10} گویشور فارسی، دست‌کم دو نفر بالای
      \N{60} سال \\
\bottomrule
\end{tabular}
\end{table}

خروجی نهایی مورد انتظار نیز شامل چهار قلم است: یک نمونه اولیه کارا، کد منبع
مستند، یک مجموعه داده مکالمه فارسی، و یافته‌های مربوط به نحوه تعامل سالمندان
فارسی‌زبان با سامانه. برای قلم چهارم، پروتکل و زیرساخت سنجش آماده شده است؛
یافته تجربی آن تنها پس از تکمیل مطالعه انسانی قابل ادعا خواهد بود و ازاین‌رو
در این نوشتار ادعایی درباره تجربه سالمندان مطرح نمی‌شود.

\section{چالش‌های اصلی}
\label{sec:challenges}

در جریان اجرای پروژه، چهار چالش تعیین‌کننده ظاهر شد که شکل نهایی کار را
تغییر داد.

\paragraph{نبود یک مسیر آموزشی رسمی برای فارسی.}
مدل \en{Mini-Omni2}~\cite{xie2024miniomni2} که در تعریف پروژه به‌عنوان نمونه
مدل پایه نام برده شده است، و همچنین
\en{LLaMA-Omni2}~\cite{fang2025llamaomni2}، مخزن‌هایی برای استنتاج منتشر
کرده‌اند اما آموزش‌دهنده کامل و رسمی گفتار‌به‌گفتار ارائه نمی‌دهند. در میان
مدل‌های بررسی‌شده، تنها خانواده \en{Moshi} یک آموزش‌دهنده رسمی تطبیق کم‌پارامتر
منتشر کرده بود~\cite{moshifinetune2024}؛ به همین دلیل مسیر مستقیم این پروژه بر
پایه \en{Moshika}، نسخه هفت‌میلیارد‌پارامتری این خانواده، بنا شد.

\paragraph{نبود داده مکالمه‌ای برچسب‌دار فارسی.}
هیچ پیکره عمومی فارسی با برچسب وقفه و هم‌پوشانی در اختیار نبود. مجموعه داده
این پروژه ناچار از یک آرشیو داخلی محتوای گفتاری فارسی و با یک زنجیره خودکار
شامل بازسازی پنجره، تفکیک گوینده و استخراج جفت پاسخ ساخته شد. پیامد این انتخاب
آن است که برچسب‌های تعاملی حاصل «خودکار»اند و نه تأییدشده توسط انسان؛
این مرز در سراسر نوشتار حفظ شده است.

\paragraph{فاصله میان کاهش زیان و تولید معتبر.}
مهم‌ترین یافته منفی پروژه همین است. تطبیق کم‌پارامتر مدل مستقیم، در شش نسل
آزمایشی، زیان «تحمیل معلم»\پاورقی{teacher forcing} را به‌طور پیوسته کاهش
داد، اما هیچ‌گاه به تولید «خودرگرسیو»\پاورقی{autoregressive} معتبر گفتار
فارسی نرسید. این مشاهده نشان می‌دهد که کاهش زیان آفلاین به‌تنهایی معیار
کافی برای انتخاب مدل گفتار‌به‌گفتار نیست.

\paragraph{مرز میان شواهد مهندسی و ادعای علمی.}
سنجه‌های خودکار در دسترس بودند، اما ارزیابی انسانی، سخت‌افزار هدف و
برچسب‌گذاری مستقل وقفه در مهلت پروژه فراهم نشد. تصمیم روشمند پروژه این بود که
هر سنجه با «کلاس شواهد» خود برچسب‌گذاری شود و هیچ تقریب خودکاری به
جای شواهد رسمی ارائه نگردد. این تصمیم در بخش~\ref{sec:evidence-first} به یک
روش کاری تبدیل شده است.

\section{راهبرد دومسیره}
\label{sec:two-track}

برای کاهش ریسک اجرایی، پروژه دو مسیر را به‌طور موازی پیش برد:

\begin{enumerate}
\item \مهم{مسیر مستقیم.} تطبیق کم‌پارامتر مدل \en{Moshika} با داده مکالمه
      فارسی، به این امید که یک مدل واحد گفتار‌به‌گفتار فارسی به دست آید. این
      مسیر پاسخ مستقیم به روح تعریف پروژه است.
\item \مهم{مسیر آبشاری.} یک زنجیره ماژولار شامل بازشناسی گفتار فارسی، تولید
      پاسخ متنی با یک مدل زبانی بزرگ و تبدیل متن به گفتار، که کنترل
      تمام‌دوطرفه و منطق وقفه به‌صورت جداگانه روی آن سوار می‌شود. این مسیر
      تضمین می‌کند که پروژه یک سامانه کارکردی قابل نمایش داشته باشد.
\end{enumerate}

هر دو مسیر از یک مجموعه داده مشترک استفاده کردند. نتیجه نهایی، همان‌طور که در
فصل~\ref{ch:results} گزارش می‌شود، نامتقارن بود: مسیر آبشاری از همه شروط
خودکار ازپیش‌تعریف‌شده عبور کرد و به‌عنوان نمونه اولیه تحویل‌شده انتخاب شد،
در حالی که مسیر مستقیم به‌عنوان یک نتیجه منفی مستند، همراه با یک سیگنال
یادگیری مثبت، گزارش شده است.

\section{دستاوردهای این پروژه}
\label{sec:contributions}

دستاوردهای این پروژه را می‌توان در پنج مورد خلاصه کرد.

\begin{enumerate}
\item \مهم{یک نمونه اولیه کارکردی گفتار‌به‌گفتار فارسی با کنترل
      تمام‌دوطرفه.} سامانه، جریان پیوسته صوت \en{PCM16} را از مرورگر دریافت
      می‌کند، با یک آشکارساز غلتان بر پایه انرژی، زیر‌و‌بمی و \en{MFCC} درباره
      وقفه تصمیم می‌گیرد، منبع صوت در حال پخش را متوقف می‌کند و پیش‌بافر
      \N{450} میلی‌ثانیه‌ای میکروفون را به نوبت بعدی منتقل می‌کند. در یک پنل
      ثابت و ازپیش‌قفل‌شده اعتبارسنجی، زنجیره مکانیکی با پاسخ‌گوی
      \en{Qwen2.5-0.5B} در هر \N{9} ردیف رونویسی فارسی، پاسخ فارسی و صوت
      غیرساکت تولید کرد و هیچ‌گاه به پاسخ جایگزین قاعده‌محور بازنگشت.
      پاسخ‌گوی متنی سامانه تحویل‌شده سپس به \en{Qwen3-4B} با پرامپت \en{v2}
      ارتقا یافت؛ پنل مکانیکی نه‌ردیفی پس از این ارتقا دوباره اجرا نشد.

\item \مهم{یک مجموعه داده مکالمه فارسی مشتق‌شده و ممیزی‌شده.} از یک موجودی
      \N{775.887} ساعتی، مجموعه‌ای شامل \N{6,754} جفت «گفتار کاربر / پاسخ» با
      \N{108.584} ساعت صوت استریو ساخته شد که درون بازه \N{100} تا \N{200}
      ساعت مورد نظر تعریف پروژه قرار دارد. تفکیک آموزش، اعتبارسنجی و آزمون بر
      پایه شناسه نشست منبع انجام شد و ممیزی خودکار، نشت گروهی صفر را تأیید
      کرد.

\item \مهم{یک نتیجه مثبت خودکار برای کیفیت معنایی پاسخ.} پس از عبور از یک
      مرحله صلاحیت هفت‌ردیفی، آزمون نهایی \N{40} ردیفی گروه‌ناهمپوشان از هر
      هفت شرط خودکار ازپیش‌تعریف‌شده عبور کرد: میانگین ارتباط پاسخ از
      \N{1.425} به \N{3.150} و میانگین انسجام به \N{3.325} رسید، با نرخ برد
      زوجی \N{67.5} درصد و اعتبار کامل هر \N{240} فراخوانی داور. داور و
      نامزد هم‌خانواده‌اند و این سنجه جایگزین ارزیابی انسانی نیست.

\item \مهم{یک نتیجه منفی مستند و قابل استناد برای مسیر مستقیم.} شش نسل
      آزمایشی، هر یک با تغییر یک عامل قابل ابطال، طراحی و اجرا شد. آخرین نسل
      کاهش \N{32.20} درصدی زیان متنی درون‌نمونه را نشان داد اما هر چهار
      چک‌پوینت آن در آزمون اجرای خودرگرسیو نه‌ردیفی \N{0} از \N{9} گرفتند.
      این تفکیک میان «ظرفیت یادگیری» و «توانایی تولید» یافته علمی اصلی این
      نوشتار است.

\item \مهم{یک روش کاری شواهد‌محور.} همه پروتکل‌های ارزیابی، شاخص ردیف‌ها و
      آستانه‌های پذیرش پیش از اجرا قفل و ثبت شدند، مصنوعات نتیجه یک‌بار‌نوشت
      هستند، و یک تجمیع‌کننده «شکست‌بسته»\پاورقی{fail-closed} مانع از
      ارتقای تقریب‌های خودکار به شواهد رسمی می‌شود.
\end{enumerate}

برای چهار نیازمندی هنوز شواهد رسمی کامل گردآوری نشده است: تأخیر سرتاسری
\N{500} میلی‌ثانیه، عبور رسمی تشخیص وقفه از \N{80} درصد با برچسب مستقل،
ارزیابی روی پردازنده گرافیکی فیزیکی هدف، و مطالعه انسانی. دامنه دقیق
نتایج موجود و این محدودیت‌ها در فصل~\ref{ch:discussion} گزارش شده‌اند.

\section{ساختار این نوشتار}
\label{sec:structure}

این نوشتار در هفت فصل سازمان یافته است.

فصل~\ref{ch:prelim} مفاهیم اولیه مورد نیاز را معرفی می‌کند: نمایش سیگنال گفتار
و ویژگی‌های آوایی، بازشناسی گفتار و تبدیل متن به گفتار، مدل‌های زبانی بزرگ و
تطبیق کم‌پارامتر، رمزینه‌های صوتی عصبی و معماری مدل‌های گفتار‌به‌گفتار، و
سنجه‌ها و مفاهیم آماری به‌کاررفته در ارزیابی.

فصل~\ref{ch:literature} کارهای پیشین را بررسی می‌کند و جایگاه این پروژه را در
میان سامانه‌های آبشاری و مدل‌های سرتاسری گفت‌وگوی گفتاری، منابع گفتاری فارسی و
پیشینه تشخیص وقفه مشخص می‌سازد.

فصل~\ref{ch:method} روش پیشنهادی را شرح می‌دهد: زنجیره ساخت مجموعه داده،
معماری سامانه آبشاری، ماژول تشخیص وقفه و بستر انتقال تمام‌دوطرفه، طراحی
آزمایش‌های مسیر مستقیم، و روش کاری شواهد‌محور.

فصل~\ref{ch:results} نتایج را با جداول و شکل‌های مربوط ارائه می‌کند و برای هر
نتیجه، مرز ادعای مجاز آن را می‌نویسد.

فصل~\ref{ch:discussion} به تفسیر نتایج می‌پردازد، به‌ویژه فاصله میان کاهش زیان
و تولید معتبر، و سپس محدودیت‌ها و ملاحظات اخلاقی و حقوقی را برمی‌شمارد.

فصل~\ref{ch:conclusions} جمع‌بندی و پیشنهادهای کار آینده را ارائه می‌کند. در
پیوست نیز فرمان‌های بازتولید، فراپارامترهای دقیق آزمایش‌ها و نقشه مصنوعات
شواهد آمده است.
